\section{Isogeny, transfer and structure routes}\label{sec:transfer}

A recurring hope was that some other representation of the same
instance---an isogenous curve, a curve at another level of the isogeny
volcano, a Jacobian, a Weil descent---would make index calculus cheaper.
These routes produced the programme's clearest negatives, and one lemma
that explains most of them.

\begin{lemma}[transport conservation,
\art{docs/ic/RESEARCH_TRANSFER_AWARE_IC_20261010.md} \S1]\label{lem:conservation}
Let $\varphi:E\to E'$ be an $\FF_q$-rational isogeny (or any homomorphism
of rational points), $\kappa=\ker\varphi\cap E(\FF_q)$, $R\in E(\FF_q)$,
$\Fb'\subset\varphi(E(\FF_q))$ a factor base on the codomain and
$\Fb=\varphi^{-1}(\Fb')\cap E(\FF_q)$. Then the number of ordered
$m$-decompositions of $R$ over $\Fb$ is $\abs\kappa^{m-1}$ times the number
of ordered $m$-decompositions of $\varphi(R)$ over $\Fb'$.
\end{lemma}
\begin{proof}
Images of decompositions of $R$ are decompositions of $\varphi(R)$.
Conversely, for $\sum Q_i=\varphi(R)$ choose rational preimages $P_i$ of
the $Q_i$; then $\sum P_i-R\in\kappa$, and of the $\abs\kappa^m$ choices
exactly $\abs\kappa^{m-1}$ sum to $R$ (translate $P_1$ by the kernel
element).
\end{proof}

So a homomorphic transfer changes nothing but three curve-specific inputs:
the rational kernel symmetry $\abs\kappa^{m-1}$ (bounded by the rational
torsion; trivial on prime-order curves), the factor-base geometry
available on the destination, and the cost of arithmetic there. Sequences
of isogenies compose to one homomorphism; a loop is an endomorphism and
acts on the subgroup as a known scalar, useful only as a known-log
multiplier. Multi-curve factor bases glued by isogenies are single-curve
bases in disguise.

\paragraph{Conductor gaps in the ECC2K-130 class.} Jao, Miller and
Venkatesan~\cite{jmv2005} prove random self-reducibility among curves with
the same endomorphism ring, leaving open whether curves whose conductors
differ by a large prime could differ in hardness. The ECC2K-130 class has
conductor gap 263; we computed its structure
(\art{docs/ic/RESEARCH_ISOGENY_CONDUCTOR_GAP_20261007.md}) and ran
several rounds of experiments. Native counted-$S$ index calculus on two fixed CM
pairs gave floor/crater ratios $0.975$ (99\% $[0.603,1.451]$) at
$\ell=8209$ and $1.127$ $[0.922,1.342]$ at $\ell=131101$: no level effect
(192 verified records, 96 exact replays;
\art{research/ecbench_conductor_gap_20261012/}). Two-torsion
symmetrisation is a quotient isogeny: at a matched base of abscissae it
cut F4 row-reduction work $4.1$--$4.3\times$ ($\abs H=2$) and
$17.6$--$19.6\times$ ($\abs H=4$) over 1{,}128 systems---but a correction
the next day showed that at a matched number of \emph{columns} the saving
disappears; it was the cost of a redundant orbit base. A combined Weil
descent used $234$--$529\times$ the F4 operations of the blind system
(\art{research/isogeny_conductor_gap_round3_20261011/}).

\paragraph{Fibre-aware index calculus.} Isogeny fibres are torsors whose
arithmetic is controlled by one cohomology class
(\art{RESEARCH_ISOGENY_FIBERS_ECDLP.md}); relation generation is
isogeny-equivariant at zero lifting cost, and a fibre-closed pair table was
the best lever of the conductor-gap rounds ($12.8\times$ rho in its best
cell). Measured against matched rho across sizes at fixed cofactor, the
envelope IC/rho ratio grows from $8.4$--$8.8$ to $35.1$--$37.0$ ($h=4$) and
from $13.4$--$20.1$ to $60.4$--$63.1$ ($h=72$), as $\#E^{1/6}$: the gap
widens, it does not close (816 audited records;
\art{research/fiber_ic_scaling_20261011/}).

\paragraph{Other structure.} Symmetrised $u$-frame bases saved
$0.84\times$ at $n=17$ and $0.76$--$0.77\times$ at $n=23$, but cost
$4.01\times$ at $n=31$ (\art{research/koblitz_symmetrised_e2e_20260927/}).
GHS and hyperelliptic covers: ECC2K-130 does map to hyperelliptic curves,
but the genus that would pay is missing. In odd characteristic the cubic
GHS transfer from $E_0(\FF_{q^3})$ to a genus-3 Jacobian now runs end to end
with double-large-prime index calculus, at $q^{1.380}$ against rho's
$q^{3/2}$, still $1.7$--$2.6\times$ the negation-map rho model at
$q=2^{20}$ (\S\ref{sec:ladder}). The
isogeny-loop GLV/GLS work produced a measured $1.09$--$1.15\times$ faster
four-dimensional scalar multiplication (\art{research/isogeny_loop_glv_20261012/})
and a $\sqrt{\text{\'elu}}$ kernel cost exponent $0.822$ $[0.805,0.838]$;
these speed up curve arithmetic, not the discrete logarithm.

\paragraph{Named obstructions.} Blocked requirements are kept open as
problem records (\art{problems/README.md}), among them the complete-cost
calibration of root-finding IC
(\art{problems/rootfinding-complete-cost-calibration.md}), with exit
criteria rather than abandoned.
